% 本文件由 LaTeX 排版 Agent 根据有效 translation.json 自动生成。
% author=Augustine of Hippo; work_id=1312; title=论说谎

\chapter{论说谎}

% structure: p0001 source=h1 target=omitted reason=page-title-already-rendered-by-parent

\WorkTitle{De mendacio}\par

本书据其在\WorkTitle{订正录}中的位置，应写于公元395年左右，因为它是第一卷所列的最后一部作品，该卷收录了他做主教前所写的作品。有些版本认为它是写给\Person{康森提乌斯}的，但手稿并非如此。手稿可能正确，因为他后来多年后针对普里西利安派案件答复\Person{康森提乌斯}的询问，写了另一部关于该主题的作品。——\EditorialNote{本笃版编者}\par

\Emph{选自\WorkTitle{订正录}第一卷，最后一章：}\par

我也写了一本\WorkTitle{论说谎}的书，尽管花了些力气才能理解，但其中包含许多对锻炼心智有用的内容，更多对道德有益的内容，在于灌输热爱说实话。我也曾打算将其从我的作品中删除，因为它在我看来晦涩、错综复杂且令人烦恼；因此我未曾将其公之于众。后来我又写了另一本书，题为\WorkTitle{驳说谎}，我更坚决地决定并命令前一本应不再存在；然而这并未实现。因此，在我作品的这次修订中，既然我发现它仍然存在，我便命令保留它；主要是因为其中有一些必要的、在另一本中没有的内容。为什么另一本标题为\WorkTitle{驳说谎}，而这一本为\WorkTitle{论说谎}，原因在于：前者通篇公开攻击说谎，而后者有很大一部分是讨论赞成与反对的问题。然而，两者都指向同一目标。本书开头是这样写的：\Quote{\Emph{论说谎是一个大问题。}}\par

1. 关于说谎，有一个大问题，它常常在我们日常事务中出现，并给我们带来很多麻烦，以免我们或鲁莽地把非谎言称为谎言，或断定有时说谎是对的，即一种诚实、善意、爱德的谎言。我们将通过同寻求者一同寻求来费力地讨论这个问题：是否有什么好的结果，我们不必自己断言，细心的读者会从讨论过程中充分领会。的确，它充满了黑暗的角落，有许多洞穴般的曲折，常常躲避寻求者的热忱；以至于有时所发现的似乎从手中滑落，瞬间又重现，然后又再度消失。最后，追捕将更确定地逼近，并追上我们的判决。其中若有任何错误，但因真理是使人脱离一切错误的，假话是使人陷入一切错误的，我想，一个人绝不会比因过于爱真理、过于拒绝假话而犯错误更安全。因为那些大加指责的人说这太过分了，而也许真理最终会说，这还不够。但读者啊，在你读完整个作品之前，你最好不要指责；这样你就会少一些指责。不要追求辞藻：我们专注于事物，以及迅速处理一件与我们日常生活密切相关的事务，因此在措辞上我们很少花费功夫，几乎毫无功夫。\par

2. 因此，暂且搁置玩笑——它们从未被视为谎言，因为它们在语调中，并在开玩笑者本人的情绪中带着最明显的标志，表明他无意欺骗，尽管他所说的并不真实；至于这类言谈是否适合成全之人使用，那是另一个问题，我们目前不打算澄清——但搁置玩笑之后，首先要注意的是，不可将没有说谎的人视为说谎者。\par

3. 为此，我们须看清什么是谎言。因为并非每一个说出假话的人都在说谎，如果他相信或认为他所说的是真的。在相信与认为之间有这样的区别：有时，相信者感到他并不知道他所相信的（尽管他可能明知自己不知道某事，却仍对此毫不怀疑，如果他坚信的话）；而认为者则以为他知道他所不知道的。因此，凡是说出他心中所持为相信或认为的东西，即便是假的，他也没有说谎。因为他这是对自己言论信实的责任，即他由此表达他心中所持的，并以其表达的方式呈现出来。尽管他没有说谎，但若他相信了不该相信的，或以为他知道他所不知道的（即使那是真的），他也不是没有过错的：因为他把未知的事当作已知的。所以，说谎的人，是心中有一个想法，而用言语或任何符号表达另一个想法的人。因此，说谎者的心被称为双重的；即有两种思想：一种是他知道或认为是真而不表达的东西；另一种是他知道或认为是假而用以代替的东西。由此导致，他可能说出假话而不撒谎，如果他以为事情就像他所说的那样，尽管并非如此；他也可能说出真话而撒谎，如果他以为那是假的却当作真的说出，尽管实际上正如他所说。因为根据他内心的感觉，而不是根据事物本身的真假，来判断他是否说谎。因此，说出假话当作真话，而他以为那是真话的人，可被称为错误和鲁莽；但不能正确地说他说谎，因为他说话时心不二，也不愿欺骗，而是被欺骗了。但说谎者的过错，在于表达内心时存有欺骗的意愿；无论他是否真的欺骗（不论是因为他说的假话被人相信，还是因为他虽存欺骗之意但未被相信，或他说出自己以为假、实际为真的话而被人相信，这种情况下他虽存欺骗之意却未欺骗），除非在某种程度上他被人以为知道或认为正如他所说，那他就成了欺骗。\par

4. 但可能有一个非常微妙的问题：完全没有欺骗的意愿时，说谎是否完全不存在？试设想以下情形：某人说出一句他自己认为虚假的话，其理由是以为他人不会相信他，目的正是通过这种方式违背自己的信誉以劝阻听他说话的人——因为他觉察到那人本就不愿相信他。因为这里有人存心不欺骗却说了谎，如果说谎就是说出任何与你所知或所信不符的话。但倘若只有说出的话带有欺骗的意愿才算说谎，那么这个人就不算说谎：他明知或自以为说的是假话，但故意这样说，目的却是让听他话的人因不相信他而不被欺骗，因为说话者知道或认为对方不会信他。由此，如果可能有人故意说假话以免听者受骗，那么相反的情况也存在：有人故意说真话以便欺骗。因为假如一个人因为觉察到自己不被相信而决定说真话，那么他正是为了欺骗而说真话：因为他知道或认为，由于他的话出自他口，听者会认为它是假的。因此，他故意说真话以便让人以为是假话，即是为了欺骗而说真话。于是可以追问：谁更算说谎？是说假话以避免欺骗的人，还是说真话以便欺骗的人？前者明知或自以为说的是假话，后者明知或自以为说的是真话？因为我们已说过，一个人若不知道他说的是假话，且以为它是真的，就不算说谎；而即便说出真话，若他以为它是假的，反倒算说谎——因为评判的依据是他们内心的意念。因此，关于我们提出的这两种人，存在不小的疑问。第一种人明知或自认为说的是假话，且故意如此以免欺骗：例如，他知道某条路有强盗埋伏，担心某个他所关心的人会走那条路，而他知道那人并不信任他，于是他对那人说那条路没有强盗，目的正是让那人因为他说没有而以为一定有强盗，从而不会走那条路——因为那人决心不信他，视他为说谎者。第二种人明知或自认为说的是真话，却故意借此欺骗：例如，他对一个不信他的人说，那条他确实知道有强盗的路上有强盗，目的正是让听者因为以为他的话是假的而特意走那条路，从而落入强盗之手。那么，这两种人谁在说谎？是选择说假话以避免欺骗的那位？还是选择说真话以便欺骗的那位？前者说假话旨在让听者相信真相？后者说真话旨在让听者相信谎言？或许两人都说了谎？前者因为他想说出假话，后者因为他想欺骗？或者恰恰相反，两人都没说谎？前者因为他没有欺骗的意愿，后者因为他有说真话的意愿？因为现在要问的不是谁犯了罪，而是谁说了谎：确实，我们马上就能看到后者犯了罪，因为他通过说真话导致他人落入强盗之手；而前者非但没有犯罪，甚至行了善，因为他说假话使那人避免了毁灭。但这些例子也可以反过来：前者可能希望他所不愿欺骗的人遭受更严重的痛苦，因为有许多情况，人们因知道某些真事而自取灭亡，如果那些事本应不让他们知道；后者可能希望他所要欺骗的人获得某种便利，因为有些实例中，有人若得知亲人真的遭遇不幸就会自尽，而因认为那是假消息而保全了自己——因此受欺骗对他们有利，正如得知真相有时对他人有害。所以问题不在于前者说假话以免欺骗或后者说真话以便欺骗是出于行善或作恶的意图；而是撇开听话者的便利或不便，仅就真假本身而言，追问两人是否都说了谎，还是都没有说谎。因为如果说谎是怀有说出假话的意愿而发出的言语，那么那个故意要说假话且说了出来的人就算说谎，尽管他意在避免欺骗。但如果说谎是任何带有欺骗意愿的言语，那么前者不算说谎，后者才算——他甚至在说真话时也意欲欺骗。而如果说谎是带有任何虚假意愿的言语，那么两人都说了谎：因为前者愿意自己的言语是假的，后者则愿意别人对他的真话产生虚假的相信。进一步，如果说谎是怀有说出假话以便欺骗的意愿之人的言语，那么两人都没有说谎：因为前者在说假话时希望别人相信真事，后者说真话是为了让别人相信假事。因此，我们若在必要时说出自己所知为真或理应相信的真话，并且希望别人相信我们所说的话，就能免于一切鲁莽和一切谎言。然而，如果我们要么把假事当作真事，要么把未知的事当作已知，要么相信不应该相信的事，要么在不必要时说出来，但除此之外别无目的，只希望别人相信我们所说的话——那么，我们固然未能免于鲁莽错误的过失，却免于一切谎言。因为只要内心良心清白，知道自己所说的话是它认为真、相信真或以为真的，并且不希望别人相信除自己所说的之外的任何东西，那么任何一种定义都无需害怕。\par

5. 然而，谎言是否有时有用，这是一个更大更值得关切的问题。如上所述，当一个人没有欺骗的意愿，甚至专门设法使听者不受欺骗，尽管他希望自己所说的事情是假的，但目的却在于使人相信一个真理，这是否算谎言？再者，当一个人为了欺骗而故意说出真理，这是否算谎言？这些是可以怀疑的。但没有人怀疑，当一个人为了欺骗而故意说出假话时，这就是谎言。因此，出于欺骗意愿的虚假表达显然是谎言。但这是否是唯一的谎言，则是另一个问题。同时，就这种大家公认的谎言而言，让我们探讨一下，为了欺骗而故意说出假话是否有时有用。持肯定意见的人列举事例来支持他们的观点：他们举出\Person{撒辣}的例子，她笑了之后，却对天使否认自己笑过；\Person{雅各伯}受到父亲询问，回答说他是长子\Person{厄撒乌}；同样，埃及的接生妇为了拯救希伯来婴儿免于出生时被杀，说了谎，而且得到了天主的赞许和奖赏。他们挑选了许多诸如此类的例子，都是关于那些你不敢责备的人所说的谎言，因此你必须承认，说谎有时不仅是无可责备的，甚至是值得称赞的。他们还举出一个例子来说服那些不仅虔心研究圣经的人，而且说服所有人和常识：假设有一个人逃到你这里，通过你的谎言，他可能免于死亡，难道你就不说这个谎吗？如果一个病人问了一个他不宜知道的问题，即使你不回答也会使他更加痛苦，你敢为了保全他的性命而说出真相，或者宁愿保持沉默，而不愿通过一个德行而慈悲的谎言来帮助他虚弱的身体吗？通过这些和类似的论证，他们认为自己充分证明了，如果行善的需要要求，我们有时可以说谎。\par

6. 另一方面，那些说我们绝不可说谎的人，更强烈地辩护，首先引用了天主的权威，因为在十诫中明确写着：\BibleQuote{你不可作假见证}，这个总括性术语涵盖了所有谎言。因为凡说话的人都在为自己的心作见证。但恐怕有人会争辩说，并非每个谎言都称为假见证，那么他该如何回答经上所写的：\BibleQuote{说谎的口杀害灵魂}：又恐怕有人以为这话可以理解为某些说谎者例外，让他读另一处经文：\BibleQuote{一切说谎言的人，你必消灭。}所以主亲口说：\BibleQuote{你们的话，是就说是，非就说非；多余的话，都出于邪恶}。因此，宗徒在命令脱去旧人（旧人名下包含一切罪恶）时，紧接着说：\BibleQuote{所以你们应戒绝谎言，彼此言谈实话}。\par

7. 他们也不承认自己被那些从旧约中引用来作为说谎例证的经文所慑服。因为在旧约中，每件事都可能具有象征意义，尽管它确实发生过。当事物或言论以象征的方式呈现时，它并非谎言。因为每一言说都应指向它所言说的对象。但当任何事物或言论以象征的方式被行出或说出时，它向那些为着理解它而被提出的人言说它所象征的意义。由此，我们可以相信，那些被列为权威的先知时代的人物，在他们所记载的一切中，他们的行动和言语都是先知性的。同样，在他们生活中所有被同一先知之神认为值得记载于书的事件中，也都含有先知性的意义。至于那些接生妇，他们确实不能说这些妇女是通过先知之神，为了预示未来的真理，而向\Person{法郎}说一事当另一事（尽管那神确实藉着她们预示了什么，但她们并不知道自己身上所发生的事）。然而，他们说，这些妇女按照她们的程度得到了天主的赞许和奖赏。因为一个习惯为了害人而说谎的人，如果开始为了行善而说谎，那么这个人已经取得了很大的进步。但一件事本身被认为是可称赞的，与另一件事因比更坏的事而被优先选择，是不同的。一个人身体健康时所表达的祝贺，与一个病人正在好转时所表达的祝福，是不同性质的。在圣经中，甚至\Place{索多玛}在与\Place{以色列}子民的罪恶相比较时，也被称为义人。他们用这个规则来评判所有从旧约经书中引用的、未被责备或不能责备的说谎事例：要么它们被赞许为向着改善和更好事物的希望所取得的进步，要么由于某种隐藏的意义，它们根本就不是谎言。\par

8. 为此，从新约各书卷中——除了吾主所使用的\OriginalTerm{预象性预示}{figurative pre-significations}——若你考量圣人们的生活与品行、他们的言行，便不能举出任何应引发人效法撒谎之同类事例。因为\Person{伯多禄}与\Person{巴尔纳伯}的\OriginalTerm{假装}{simulation}不单被记载，更受责备与纠正。因为并非如一些人所想的，\Person{保禄}宗徒出于同样的假装而行割损；他要么为\Person{弟茂德}行了割损，要么亲自按犹太礼规举行某些仪式；他这样做，是出于他心中的那种\OriginalTerm{自由}{liberty}，他宣讲外邦人不因\OriginalTerm{割损}{circumcision}而更优，犹太人不因割损而更劣。因此他判定，前者不应受犹太习俗束缚，犹太人也不应被劝阻遵守祖传习俗。因此有他的话：\BibleQuote{有人蒙召时已受割损么？就不要使割损消失。有人蒙召时未受割损么？就不要受割损。割损算不得什么，不割损也算不得什么，只在于遵守天主的诫命。每人该在各自蒙召的身份上安分守己。}人如何在受割损后成为未受割损者？但他却说：不要这样做；他不要人如此生活，仿佛他已成了未受割损者，亦即仿佛他用肉皮重新覆盖了裸露的部位，不再作犹太人；如同他在另一处所说：\BibleQuote{你们的割损反倒成了不割损。}宗徒说这话，并非要强迫那些人留在不割损中，或强迫犹太人遵守祖传习俗；而是叫这两方都不应被迫改从对方的习俗；每人应有能力持守自己的习俗，而非有此必要。因为若一个犹太人希望从犹太规条中退后，只要不搅扰任何人，也不会被宗徒禁止，因为他劝他们持守其中的目的，是免得犹太人因多余之事困扰，而妨碍他们达至救恩所需之事。同样，若一个外邦人希望受割损，为表示他并不厌恶割损以为有害，而是视之为无关紧要的标记，其用处已随时日过时，他也不会被禁止；因为从前的救恩已不再来自割损，但也不致因此而有毁灭须惧怕。正因如此，\Person{弟茂德}虽蒙召时未受割损，但因他的母亲是犹太人，为争取亲属，他须向他们表明自己在基督徒纪律中并未学到憎恶\OriginalTerm{旧约圣事}{sacraments of the old Law}，遂由宗徒行了割损；借此他们可向犹太人证明，外邦人不领受这些圣事，并非因为它们是邪恶的、祖辈曾有害地遵守它们，而是因为在\OriginalTerm{伟大圣事}{that so great Sacrament}来临之后，它们不再是救恩所必需的；这伟大圣事是那整个古老圣经在\OriginalTerm{先知预象}{prophetical prefigurations}中长期孕育而生的。他本来也要给\Person{弟铎}行割损，当犹太人催迫时，但那些私引进来的假弟兄希望这样行，为能散布关于\Person{保禄}自己的信息，作为他屈服于他们宣讲真理的记号；他们说，\OriginalTerm{福音救恩}{Gospel salvation}的希望在于肉身割损和那类规条，没有这些，基督于人毫无益处；反之，基督绝不会有益于那些因认为割损中有救恩而受割损的人；因此有那话：\BibleQuote{看，我保禄告诉你们：若你们受割损，基督对你们就毫无益处。}因此，\Person{保禄}出于这种自由，遵守了他祖辈的规条，但带着唯一的谨慎和明确声明：人不应以为没有这些规条就没有基督徒的救恩。然而，\Person{伯多禄}因他的作为仿佛救恩在于犹太教，便是在强迫外邦人\OriginalTerm{犹太化}{judaize}；正如\Person{保禄}的话所表明的：\BibleQuote{你为什么强迫外邦人按犹太人的方式生活呢？}因为那些人若不见他遵守它们，仿佛没有它们便没有救恩，就不会受强迫。因此，\Person{伯多禄}的假装不能与\Person{保禄}的自由相提并论。我们固然因\Person{伯多禄}乐意接受纠正而应爱他，但绝不可借\Person{保禄}的权威来支持撒谎；\Person{保禄}曾在众人面前将\Person{伯多禄}引回正途，免得外邦人因他而被迫犹太化；同时他为自己宣讲作证：他虽然因不愿将祖传规条强加于外邦人而被视为敌视传统，但他本人并不轻视按祖传习俗举行这些仪式，并充分表明在基督来临后这些仍然保留；因此，它们对犹太人并非有害，对外邦人亦非必要，且从此对人类中的任何人都不再是救恩的手段。\par

9. 但如果不能从古经中举出任何为谎言辩护的权威——或因认为以寓意方式理解的行为或言语并非谎言，或因善人不必效法恶人改过之初在比较中受称赞的行为；也不能从新约中举出，因我们应效法的是伯多禄的纠正而非他的伪装，他的眼泪而非他的否认——那么，对于那些取自日常生活的例子，他们更确信不可信赖这些。首先，他们用圣经的多项证据教导说，谎言是不义，尤其是那段经文：\Quote{主，你恼恨一切作不义的人，你必消灭说虚谎的人。} 因为圣经惯常在下文解释上文，故\Quote{不义}是广义词，而\Quote{虚谎}是特指的那种不义；或者，若他们认为两者有区别，则\Quote{虚谎}比\Quote{不义}更严重，正如\Quote{你必消灭}比\Quote{你恼恨}更重。因为天主或许较温和地恼恨某人而不消灭他，但祂所消灭的人，祂恼恨得更加厉害，因祂惩罚得更严厉。祂恼恨一切作不义的人，但祂也消灭一切说虚谎的人。这一点既定，那么，主张此说的人中，谁会被那些例子所打动呢？比如有人说，假设有一个人向你寻求庇护，而你可以用你的谎言使他免于一死？因为那愚人害怕却不畏惧犯罪的那种死亡，如主在福音中所教导的，杀害的不是灵魂而是身体，因此祂嘱咐我们不要害怕那种死亡；但说谎的口杀害的不是身体而是灵魂。因为经上明明写着：\Quote{说谎的口杀害灵魂。} 那么，怎能不极度颠倒地说，为使一人得身体的生命，另一人必须遭受灵魂的死亡呢？邻人之爱以各人对自己的爱为限度。祂说：\Quote{你应当爱你的近人如你自己。} 一个人若为保全他人暂世的生命而丧失自己永恒的生命，怎能说是爱他如同自己呢？因为若为他人暂世的生命而丧失自己的暂世生命，那已是爱他超过自己，越过了健全教义的规则。那么，他若因说谎而为他人暂世的生命丧失自己永恒的生命，就更不可取了。当然，一个基督徒不会犹豫为自己的邻人永恒的生命而丧失自己的暂世生命：因为主为我们死了，这就是先例。为此祂也说：\Quote{这是我的命令：你们该彼此相爱，如同我爱了你们一样。人若为自己的朋友舍掉性命，再没有比这更大的爱情了。} 没有人会愚蠢到说主所做的不是为人的永远得救，无论祂是在做祂所嘱咐我们做的事，还是嘱咐我们做祂自己所做的事。既然说谎会丧失永恒生命，就绝不可为任何人的暂世生命而说谎。至于那些不满于别人拒绝说谎、因而杀害自己灵魂以使他人肉身长寿的人：倘若我们偷窃或奸淫能使一人免死，我们就应偷窃、行淫吗？在这类情况下，他们不能说服自己：例如，若有人拿着一根绳索，要求你顺从他的肉欲，并声称若你不答应他便上吊自杀；他们不能说服自己去为所谓救他一命而妥协。若这实属荒谬邪恶，那么为什么一个人要用谎言败坏自己的灵魂以使他人活于肉身呢？若他为了同一目的让自己的身体被败坏，在众人眼中他必犯下可憎的恶行。故此，这问题唯一应关注的是：谎言是否不义。既然上文引用的经文肯定了这一点，我们必须明白：问一个人是否应为救他人而说谎，等于问一个人是否应为救他人而行不义。但若灵魂的得救拒绝这一点——因为得救只能靠正义，并要求我们不仅将之置于他人之上，甚至置于自身暂世安全之上——那么，还剩下什么能使我们怀疑绝不应当在任何情况下说谎呢？因为不能说在暂世之物中有比身体的安全和生命更重大、更宝贵的。因此，如果连这一点也不能优于真理，那么那些认为有时可以说谎的人，能为劝人说谎而提出什么理由呢？\par

10. 论及身体的纯洁：在这方面，的确有一种看似非常可敬的理由介入，并以它的名义要求说谎；也就是说，如果可以用谎言逃脱强暴者的攻击，那么毫无疑问说这种谎是正确的。但对此很容易回答：身体的纯洁只依赖于心灵的完整；心灵一旦破损，即使身体外表看来完好，也必然败坏；因此身体的纯洁不应被视为可被违背他人意愿而夺走的暂时之物。所以，绝不可为了身体而让心灵因谎言败坏，因为只要心灵的纯洁未失，身体就保持纯洁。因暴力且没有先行的贪欲而使身体所受的，更应称为强夺而非败坏。或者，若一切强夺都是败坏，则并非每一种败坏都有羞耻，只有那些由贪欲所引起或同意的败坏才有羞耻。心灵比身体越尊贵，心灵的败坏就越严重。因此，纯洁可以在心灵中保持，因为那里只有自愿的败坏才能发生。确实，如果强暴者攻击身体，既无法用反抗之力也无法用任何计谋或谎言逃脱他，我们必须承认，纯洁不可能被别人的贪欲所侵犯。因此，既然无人怀疑心灵优于身体，我们应当将心灵的完整置于身体的完整之上，因为心灵的完整可以永远保持。那么，谁说说谎者的心灵是完整的呢？实际上，贪欲本身正当地被定义为：一种心灵的欲求，将任何暂时的好处置于永恒的好处之上。因此，没有人能证明任何时候说谎是正确的，除非他能证明通过谎言可以获得任何永恒的好处。但由于每个人远离真理的程度与远离永恒的程度成正比，说人可以通过远离真理而获得任何好处，是极其荒谬的。否则，如果存在任何不被真理包含的永恒好处，那它就不是真正的好处，因此也不是好处，因为它是虚假的。正如心灵之于身体，真理也应当被置于心灵本身之上，使心灵不仅超过身体地渴慕真理，甚至超过自身。这样，当心灵享受真理的不变性而非自身的可变性时，它就会更加完整和纯洁。如果\Person{罗特}——一位如此正义、甚至堪与天使为伍的人——为了不让男子受凌辱而将自己的女儿交给\Place{索多玛}人的贪欲，那么为了在真理中保持心灵的纯洁，我们岂不应当更加勤勉和坚定地守护它？因为心灵的纯洁确实比身体更为宝贵，甚至胜过男子的身体对于女子的身体。\par

11. 但若有人以为，为了别人而说谎是合理的，为的是让他能活下去，或不在他所深爱的事物中受冒犯，以便通过教导而获得永恒的真理——那么此人首先不明白，基于同样的理由，他将被迫犯下任何罪恶，正如上面所证明的；其次，教义的权威将被切断并完全毁灭，因为如果我们试图引领的人因我们的谎言而认为有时说谎是正确的。因为那带来救恩的教义，一部分在于当信之事，一部分在于当领悟之事；而除非先相信那些当信之事，否则无法达到那些当领悟之事；那么，一个认为有时说谎正确的人，如何能令人相信他呢？唯恐他在教导我们相信的那一刻，也因某个他认为的理由而说谎？因为谁能知道他是否在那时认为有一个善意的谎言的理由，以为用一个虚假的故事可以吓唬人、阻止人纵欲，并因此认为说谎甚至在精神事物上也是行善？一旦这种谎言被接受和认可，信仰的全部纪律就完全被颠覆；纪律被颠覆，就无法达到领悟，因为纪律正是为领受领悟而培育婴孩的；因此，全部真理的教义都被废除，让位于最放纵的虚假，如果谎言，即使是善意的，可以从任何地方找到入口进入的话。因为，要么说谎者将暂时的好处——自己或他人的——置于真理之上，还有什么比这更乖谬的呢？要么当他试图通过谎言使人适合获得真理时，他反而阻断了通往真理的道路，因为当他说谎时想要迎合他人，结果他讲真话时反而不可信赖。因此，要么我们不该相信好人，要么我们必须相信那些我们认为有时必须说谎的人，要么我们不该相信好人有时说谎：这三种情况中，第一种有害，第二种愚蠢；所以只能认为好人绝不说谎。\par

12. 至此，问题已从两方面加以考虑和讨论；但仍然不易作出判决：但我们必须进一步仔细倾听那些说\Quote{没有一样行为是如此邪恶，以至于为了躲避更大的邪恶而不该去做}的人；此外，人的行为不仅包括他们所做的，也包括他们同意别人对自己所做的。因此，如果一个基督徒因受到逼迫者的威胁，若不向偶像焚香就遭受身体玷污，而选择向偶像焚香以免于这玷污，那么他们认为有理由质问：为何不能用说谎来逃避如此可耻的羞辱？因为，他们声称，同意忍受对个人的侵犯而不向偶像焚香，这本身不是被动承受，而是一项行为；为了避免此行为，他选择了焚香。那么，他岂不更乐意选择谎言，若用谎言能阻止如此骇人的羞辱加诸圣洁的身体吗？\par

13. 在以下命题中，这些要点很可能值得探讨：这种同意是否应被视为一种行为？或者，那种没有认可的同意识是否可称为同意？或者，当有人说\Quote{忍受此苦比做那事更为有益}时，这算不算认可？以及，所提及之人是否做对了，选择焚香而不愿身体遭受侵犯？又或者，如果所提的替代方案是说谎，那么宁愿说谎而不焚香是否正确？但如果这种同意被视为一种行为，那么那些宁愿被处死也不作假见证的人就成了杀人犯，而且更糟的是，他们是自杀者。因为按此推论，为何不能说他们自杀了，因为他们选择了让自己遭受那样的待遇，以避免去做被强迫之事？或者，如果杀死他人比自己被杀更为严重，那么假若向一位殉道者提出这样的条件：如果他拒绝为基督作假见证并向魔鬼献祭，那么就在他眼前，不是另一个人，而是他自己的父亲将被处死——他的父亲恳求他不要因坚持而允许此事发生——这又该如何？难道不明显吗：如果他坚持他忠诚见证的决心，那么仅仅是那些杀死他父亲的人才是杀人犯，而他本人并非也成了弑亲者？因此，如同在这案例中，此人不会参与如此严重的罪行，因为他宁愿由他人杀死自己的父亲（即便那父亲是亵圣之人，其灵魂将被夺去受罚），也不愿通过假见证违背自己的信德；同样，在第一个案例中，类似的同意也不会使他参与那样可耻的侮辱，如果他拒绝自己作恶，无论他人因他未作恶而做了什么。因为这些迫害者说的话无非是：\Quote{你作恶，以免我们作恶。}如果情况是，我们作恶就能使他们不作恶，即便如此，我们也不应当通过自己行恶来使他们免于罪责；但事实上，当他们对这类事只字未提时，他们已经在作恶了，为何要让他们与我们同流合污，而不是让他们独自卑劣可憎？因为这种情形不应称为同意：既然我们不认可他们的行为，始终希望他们不要那样做，并尽我们所能阻止他们那样做，而当事情发生后，我们不仅不与他们同流合污，反而极其憎恶地谴责它。\par

14. 你说：\Quote{怎么？如果他们不这样做，除非你那样做，这难道不是你的作为和他们的作为一样吗？}按此推理，我们会与入室盗窃者一同入室盗窃，因为如果我不关门，他们就不会撬门；我们会与拦路抢劫者一同杀人，如果我们碰巧知道他们要杀人，因为如果我不当场杀死他们，他们就会杀死别人。或者，如果有人向我们坦白他将要弑亲，如果我们有能力在他动手前杀死他以防止或阻止他，而我们没有那样做，那么我们就是与他一同犯罪。因为可以逐字重复之前的话：\Quote{你与他一同做了这事，因为如果不是你做了那事，他就不会做这事。}依我的意愿，两者都不应作恶；但只有一件事在我能力范围内，我能确保它不被做；另一件事取决于他人，当我无法通过善意劝告打消他的意图时，我不应通过自己的恶行来阻止他的行为。因此，拒绝为他人犯罪并非认可罪人；无论哪一种恶行，那些希望两者都不发生的人都不会喜欢它们；但在与自己相关的事上，他有能力做或不做，于是他就不做；在与他人的事上，他只能意愿或不意愿，于是他谴责它。因此，当那些人提出条件说：\Quote{如果你不焚香，你就会遭受此苦}时，如果他回答：\Quote{至于我，我两者都不选，两者都憎恶，我不在任何事上同意你们}——他说出这些和类似的话，这些话因为是真实的，所以不会为他们提供任何同意或认可——那么让他承受他们可能施加的一切，他的账上记的是接受伤害，他们的账上则是犯下罪行。有人会问：\Quote{那么，他是否应当宁愿身体被侵犯，也不焚香？}如果问题是\Quote{他应当怎样}，他两者都不应当做。因为如果我说他应当做其中任何一件，我就会认可这个或那个，而我两者都反对。但如果问题是，既然无法避免两者，但能避免其一，他应当优先避免哪一件？我将回答：\Quote{宁可避免自己的罪，而不避免他人的罪；宁可是自己犯的较轻的罪，而不是他人犯的较重的罪。}因为，在保留更仔细探究的前提下，暂且承认身体侵犯比焚香更坏，但焚香是自己的行为，而身体侵犯是他人之行为，尽管发生在他身上；行为是谁的，罪就是谁的。因为虽然杀人是比偷窃更重的罪，但偷窃比被人杀害更坏。因此，如果有人向某人提出：如果他不偷窃，就会被杀——即谋杀将临到他身上——既然他无法两者都避免，他会宁愿避免那将成为他自己的罪的行为，而不是他人之罪。这后者的罪并不会因临到他身上而成为他的行为，也不会因为他可以通过犯自己的罪来避免它而成为他的行为。\par

15. 因此，这个问题的全部重点在于：是否普遍成立，即，如果你能通过自己一个较轻的罪来避免别人加于你的罪，但你却没有这样做，那么这别人的罪就不归咎于你；或者，是否所有身体上的玷污都是例外。没有人会说，一个人因被谋杀、被投入监狱、被锁链捆绑、被鞭打、被其他酷刑和痛苦折磨、被剥夺财产甚至到赤身露体、被剥夺荣誉、受各种辱骂而蒙受极大耻辱，就因此而被玷污。但是，如果有人将污秽物倒在他身上，或倒进他嘴里，或塞进他嘴里，或像对待女人一样对他进行肉体上的侵犯，那么几乎所有人都会以惊恐的目光看待他，并称他为玷污和不洁。因此必须得出结论：无论别人的罪是什么，那些加于某人时玷污他的罪除外，人不得为了自己或为了任何别人，而通过自己的罪来寻求避免这些罪，反而必须忍受并勇敢地承受；如果他不应该通过自己的罪来避免它们，那么也不应该通过谎言来避免；但是，那些加于人时使人不洁的罪，我们甚至必须通过犯罪来避免它们；因此，为了避免那种不洁而做的事，就不应该被称为罪。因为凡是做了的事，考虑到如果不做就会成为责备的理由，那么这件事就不是罪。按照同样的原则，当无可避免时，那种情况也不应被称为不洁；因为即使在那种极端情况下，遭受它的人也有可以正确做的事，即耐心地承受他无法避免的事。现在，没有人能在正确行事时被任何肉体的传染所玷污。因为在天主眼中，不洁的是每一个不义的人；因此，洁净的是每一个正义的人；即使不是在人眼中，至少在天主眼中，祂的判断无误。不但如此，即使在有避免的能力时遭受那种玷污，人也不是仅仅因接触而被玷污，而是因拒绝在他可能避免时去避免的罪。因为如果是为了避免它而做的任何事，那就不会是罪。因此，凡是为了避免它而说谎的人，并没有犯罪。\par

16. 或者，是否有些谎言也要除外，以至于宁可忍受这种玷污，也不去犯那些谎言？如果是这样，那么并非一切为了避免那种玷污而做的事就不再是罪；因为有些谎言犯下比遭受那种污秽的暴力更糟糕。例如，假设有人在追捕一个人，要玷污他的身体，而有可能通过一个谎言来掩护他；谁敢说即使在这种情况下也不应该说谎？但是，如果用来隐藏他的谎言可能损害另一个人的声誉，给他带来虚假的指控，正是那种对方正在被追捕的玷污，比如，如果对询问者说：\Quote{你去找某某人}（说出一个洁身自好、与这类恶习无关的人的姓名，）\Quote{他会替你找一个你发现更愿意的人，因为他认识并喜欢这样的人；}这样，那人就可能从他所寻找的人身上转移开。我不知道是否应该为了不让别人的身体被他不认识的情欲所玷污，而用一个谎言来损害一个人的声誉。一般来说，绝不可以为了任何人而说一个会伤害另一个人的谎言，即使那伤害比他若不说那谎言所受到的伤害要轻。因为不可以违背另一个人的意愿夺走他的面包，即使他身体健康，而面包是用来喂一个虚弱的人；也不可以违背一个无辜者的意愿用棍棒打他，以免另一个人被杀。当然，如果他们愿意，就可以这样做，因为他们如果愿意这样，就没有受到伤害。但是，是否甚至经得本人的同意，就可以用虚假的指控来损害一个人的声誉，指控他有污秽的情欲，以便情欲不从另一个人的身体上发生，这是一个大问题。而且我不知道是否容易找到一种方式，可以公正地让一个人的声誉，即使是他同意的，被虚假的情欲指控所玷污，就像一个人的身体可以违背他的意愿被情欲本身所玷污一样。\par

17. 但是，如果有人提议给那个宁愿向偶像焚香也不愿让身体遭受可憎情欲的人一个选择，即如果他想要避免那种情欲，他应该用某个谎言来损害基督的名声；他若这样做，就是极其疯狂。我不妨说：他若是为了避免别人的情欲，并且不让别人对他做他本人并无情欲的事，而用虚假的赞美基督来伪造基督的福音，那就是疯狂的；他更害怕别人玷污他的身体，而不是自己去玷污那圣化灵魂和身体的教义。因此，在宗教教义方面，以及普遍地在那些为宗教教义而说出的、用于教导和学习该教义的话语中，必须完全杜绝一切谎言。而且，在涉及这类事情时，任何人都不应认为有任何理由可以说谎，因为在这教义中，即使是为了更容易地引导人归向它，也不可以说谎。因为一旦真理的权威被打破或稍受削弱，一切都会变得不确定；除非人们相信它们是真实的，否则就不能确定地持守它们。因此，无论是谈论、辩论和宣讲永恒之事的人，还是叙述或谈论与宗教和虔敬的建造有关的暂时之事的人，都可以在适当的时候隐藏他认为适合隐藏的事；但说谎是绝对不允许的，因此也不允许通过说谎来隐藏。\par

18. 这一点从一开始就最坚定地确立了，关于其他谎言的问题就更可靠地展开了。但接下来我们也必须看到，一切伤害他人的不义谎言都必须避免：因为没有人应该受到不义对待，即使为了让另一个人免于更重的伤害而对他造成较轻的伤害，也是不允许的。那些虽然不伤害他人，但也不使任何人受益，反而对随意说谎者本身有害的谎言，也是不可允许的。实际上，这些人才是真正可以被称为说谎者的人。因为说谎与作为说谎者是有区别的。一个人可能不情愿地说了一个谎；但说谎者喜爱说谎，并在内心以说谎为乐。其次，还有一类人，他们通过说谎来取悦他人，不是为了伤害或中伤任何人（我们已经排除了那种情况），而是为了在谈话中显得愉快。这一类人与我们归为说谎者的那类人有如下区别：说谎者以说谎为乐，以欺骗本身为乐；但后者渴望通过愉快的谈话取悦他人，而且宁愿说真话来取悦他人，但当他们不容易找到令人愉快的真话时，他们宁愿说谎也不愿闭嘴。然而，他们有时很难编造一个完全虚假的故事；最常见的是在缺乏有趣内容时，将虚假与真实交织在一起。现在，这两类谎言对相信它们的人没有害处，因为他们在宗教和真理的问题上，或在自己利益和好处方面并没有受骗。他们只须判断所说的事是可能的，并相信一个他们不应贸然认为说谎的人。因为相信某人的父亲或祖父是一个好人，而事实并非如此，这有什么害处呢？或者相信他曾在\Place{波斯}服役，而他从未踏出\Place{罗马}城一步？但对说谎者本身，这些谎言造成了很大伤害：对前一类人来说，因为他们如此背离真理，以欺骗为乐；对后一类人来说，因为他们想取悦他人胜过真理。\par

19. 在毫不犹豫地谴责了这些谎言之后，接下来有一种谎言，似乎逐步上升到更好的层次，通常归因于善意和善良的人，即说谎者不仅不伤害他人，甚至还使某人受益。正是关于这类谎言，整个争论展开了：使他人受益却违背真理的人，是否伤害了自己？或者，如果唯独那种以内在不变的光照亮人心灵的才能称为真理，那么至少他违背了某种真实的事，因为尽管身体感官受骗，但他断言某事是或不是这样，而他的心灵、感官、意见或信仰都没有给他任何报告。因此，他是否因如此使他人受益而伤害自己，或者在补偿中因使他人受益而不伤害自己，这是一个大问题。如果真是这样，那么他应该通过不伤害任何人的谎言使自己受益。但这些事情是相互关联的，如果你承认这一点，必然会引出一些非常麻烦的后果。因为如果有人问，对于一个拥有过剩财富的人来说，从成千上万蒲式耳小麦中损失一蒲式耳有什么害处？而这一蒲式耳可能作为必要的食物对偷窃者有益；那么，偷窃也可以不受责备地实施，作假见证也可以无罪。还有什么比这更颠倒的呢？或者，如果另一个人偷了那一蒲式耳，你看到了，被询问时，你会为穷人诚实地说谎吗？如果你为自己的贫困这样做，你会受到责备吗？好像你有责任爱他人胜过爱自己。因此，两者都是可耻的，必须避免。\par

20. 但也许有人会认为应该加一个例外：有些诚实的谎言不仅不伤害人，反而使某些人受益，但那些掩盖和包庇罪行的谎言除外。这样，上述谎言之所以可耻，是因为它虽然不伤害人并且有益于穷人，却掩盖了偷窃；但如果它以这种不伤害人并使某些人受益的方式，而不掩盖和包庇任何罪恶，那么它在道德上并非错误。例如，假设有人在你眼前藏匿他的钱，以免被偷或遭暴力，你因此被询问时说了谎；你不伤害任何人，并且帮助了那个需要隐藏钱财的人，而且没有通过说谎掩盖罪恶。因为一个人藏匿他害怕失去的财产并不是罪。但是，如果我们因为说谎时没有掩盖任何人的罪，不伤害任何人并做了好事，就不算犯罪，那么关于谎言本身的罪，我们该怎么说呢？因为那里写着：\Quote{不可偷盗}，也写着：\Quote{不可作假见证}。既然两者分别被禁止，为什么作假见证在掩盖偷盗或其他罪恶时是有罪的，而若不掩盖任何罪恶而单独发生时，就无罪呢？而偷窃本身在任何情况下都是有罪的，其他罪恶也是如此？难道掩盖罪恶是不允许的，而犯罪却是允许的吗？\par

21. 如果这很荒谬，我们该怎么说呢？是否只有在一个人说谎编造对某人的控告、或掩盖某人的罪行、或以任何方式在审判中压迫某人时，才算是\Quote{假见证}？因为见证人似乎是法官审理案件所必需的。但若圣经中的\Quote{见证}只限于此，那么宗徒就不会说：\Quote{我们竟被发现是天主的假见证，因为我们曾为天主作证说，祂复活了基督，而其实祂没有复活。}因为这样他表明，说谎即使是虚假地赞扬某人，也是假见证。\par

或者，说谎者是否在捏造或隐瞒任何人的罪恶，或以任何方式伤害任何人时，便作了假见证？因为，如果针对人现世生命的谎言是可憎的，那么针对永生的谎言更是何等可憎？正如每一个谎言，如果发生在宗教教义中，便是如此。正是为此，宗徒称关于基督的谎言为假见证，甚至看似归于祂赞美的谎言也是如此。那么，如果一个谎言既不捏造也不隐瞒任何人的罪恶，也不是对法官问题的回答，不伤害任何人，却有益于某人，我们是否要说它既不是假见证，也不是应受谴责的谎言？\par

那么，如果一个杀人犯向基督徒寻求庇护，或者基督徒看见杀人犯躲藏的地方，并被追捕者（为了惩罚杀人者）问及此事时，他应该说谎吗？因为说谎时，他如何没有隐瞒罪恶呢？毕竟他所庇护的人犯了重罪。或者，是否因为追问的不是他的罪，而是他躲藏的地点，所以为了隐瞒罪人说谎是邪恶的，但为了隐瞒罪人本身却不是邪恶的？\Quote{确实如此，}有人说，\Quote{因为人不是在避免惩罚时犯罪，而是在做应受惩罚的事时犯罪。再者，基督信仰的纪律要求既不对任何人的悔改失望，也不阻挠任何人的悔改之路。}如果你被带到法官面前，并被问及那人躲藏的具体地点，你会准备说：\Quote{他不在那里，}（尽管你知道他在那里），或者\Quote{我不知道，也没有看见，}（尽管你知道并看见了）吗？那么你是否准备作假见证，为了不让杀人者被杀而杀害你的灵魂？或者，在法官面前你宁愿说谎，但当你被法官讯问时，你又说真话以免作假见证？那么你就是通过出卖他来杀他。的确，圣经也憎恨出卖者。或者，在法官讯问时如实回答的人不算出卖者，但未经询问主动告发别人以致其灭亡的才是出卖者？设若有一个公义无辜的人，你知他藏在哪里，被法官讯问；这人被更高权力下令处决，而讯问者只是执行法律，并非判决者。那么，为了无辜者说谎是否不算假见证，因为讯问者不是法官，只是执行者？如果立法者本人讯问你，或者任何不义的法官，为了惩罚无辜者而追查他，你该怎么办？你会作假见证还是出卖者？或者，向公正的法官主动告发一个藏匿的杀人犯算出卖者，而向不义的法官（正在追查他要杀害的无辜者）透露那个信任他的人的藏身之处就不算出卖者？或者，你会在假见证和出卖的罪之间犹豫不决？或者，通过保持沉默或声明拒绝回答，你决定避免两者？那么，为什么不在见到法官之前就这样做，也好避开谎言呢？因为，避开谎言就能避免一切假见证（无论是否每个谎言都是假见证）；但按照你们的说法，避开一切假见证并不能避开一切谎言。那么，说出\Quote{我既不告密也不说谎？}要勇敢多少，卓越多少呢？\par

塔尔加斯特教会的一位前任主教，名叫\Person{菲尔穆斯}，意志更为坚定，就这样做了。因为，当他奉皇帝之命，通过他派来的官员被问及一个向他寻求庇护、并被他尽可能小心隐藏的人时，他回答他们说，他既不能说谎，也不能出卖人；当他遭受了那么多身体的折磨时（当时皇帝尚未信奉基督），他仍坚守初衷。随即被带到皇帝面前，他的行为显得如此可敬，以至于毫不费力地为他试图拯救的人获得了赦免。还有什么行为比这更勇敢、更坚毅呢？但或许有些更胆小的人会说：\Quote{我可以准备好忍受任何折磨，甚至接受死亡，以免犯罪；但是，既然说一个不伤害任何人、不作假见证、且有益于某人的谎言并非罪过，那么自愿且无意义地忍受折磨，并在自己的健康和生命或许还有用时，将它们白白丢给愤怒的人们，便是愚蠢和严重的罪过。}我问这些人：为什么他们害怕经上记载的\BibleQuote{你不可作假见证，}却不害怕向天主所说的\BibleQuote{你必消灭一切说谎的人？}他说：\Quote{经上没有说\InnerQuote{每一个谎言}，但我理解它好像是说\InnerQuote{你必消灭一切作假见证的人}。}但那里也没有说\InnerQuote{一切假见证}。他说：\Quote{是的，但它被放在那里，与其他每方面都是邪恶的事物放在一起。}那么，那里所记载的\BibleQuote{你不可杀人}也是如此吗？如果这每方面都是邪恶的，那么那些依据所颁布的法律杀了许多人的义人，又如何能免于此罪呢？\Quote{但是，}有人反驳说，\Quote{那执行某个正当命令的人，他自己并没有杀人。}那么，我接受这些人的担忧，我仍然认为那位既不说谎也不出卖人的可敬之人，更好地理解了经上所记载的，并勇敢地将所理解的付诸实践。\par

24. 但有时会遇到这样的情况：我们没有被询问被寻找的人在哪里，也没有被迫出卖他，如果他藏得如此隐蔽，除非被出卖否则不易被发现；但我们被问，他是否在某处。如果我们知道他在那里，保持沉默就会出卖他，甚至说我们无论如何也不会说出他在不在那里：因为这样提问者就会推导出他在那里，如果他不在，那么那个既不愿说谎也不愿出卖人的人，除了说\Quote{他不在那里}之外，不会回答别的。因此，通过我们保持沉默或说这样的话，一个人被出卖了，寻找他的人只需进去（如果有权力）就能找到他；而原本他可能因我们说谎而错过找到他。因此，如果你不知道他在哪里，就没有理由隐瞒真相，而必须承认你不知道。但是，如果你知道他在哪里，无论是在被问的那个地方还是别处，当被问他在不在那里时，你绝不能说\Quote{我不会告诉你你所问的}，而必须说\Quote{我知道他在哪里，但我绝不会指出来}。因为如果对某个特定地点你不回答并声明你不会出卖，那就好比用手指指着那个地方一样：因为这会引起确定的怀疑。但如果你一开始就承认你知道他在哪里但不会说，或许询问者会转移那个地点，开始逼你透露他实际所在的地方。为了这样的诚信和仁爱，无论你勇敢地承受什么，都被视为不仅无罪，甚至值得赞扬；唯有那些如果一个人忍受了就会被称为不是勇敢地忍受，而是无耻和污秽地忍受的事情除外。因为这是关于谎言的最后一种描述，我们必须更仔细地讨论它。\par

25. 首先应避免的是那须远避的重大谎言，即在宗教教义中所说的谎言；人绝不应受任何考虑而被诱导于这种谎言。第二种，是那种不义地伤害某人的谎言：这种谎言对任何人无益，却伤害某人。第三种，是那种有利于一人而伤害另一人，但不涉及身体玷污的。第四种，是仅出于谎言和欺骗的欲望而说的，这是纯粹的谎言。第五种，是出于交谈中取悦他人的欲望而说的。这些谎言全被避免和摒弃后，接下来是第六种，它既不伤害任何人，又有助于某人；例如，当某人的钱财将被不义地夺走时，知道钱在哪里的人，无论被谁询问，都应说不知道。第七种，是不伤害任何人而有益于某些人的谎言：除非是法官讯问的情况：例如，当不愿出卖一个被搜索出来要被处死的人时，可以说谎；不仅是对一个义人和无辜者，甚至对罪人也如此；因为基督徒的规范既不应绝望于任何人的悔改，也不应阻止任何人的悔改之路。对于这两种经常引起很大争议的谎言，我们已经充分讨论过，并表明了我们的判断：即通过承担那些光荣且勇敢承受的后果，这些种类的谎言也应由勇敢、忠信和诚实的人男女所避免。第八种谎言是那种不伤害任何人，而有助于保护某人免受身体玷污，至少是上述那种玷污。因为犹太人甚至认为不洗手吃饭是一种玷污。或者，如果有人也认为这是一种玷污，但这并不至于要为了躲避而说谎。但如果这种谎言造成对某人的伤害，即使它保护一个人免于所有人都憎恶和厌恶的那种不洁；那么，只要谎言所造成的伤害不在于我们目前所关心的那种不洁，这种谎言是否可以说，是另一个问题：因为这里的问题不再是关于说谎，而是是否应当对任何人造成伤害，即使不是通过谎言，以使上述玷污可以从另一个人身上避开。我绝不认为应当如此：尽管所提出的情况是最轻微的冤屈，如我上面提到的，关于一升小麦；而且尽管很令人困惑的是，我们是否有责任不对任何人造成甚至这样的伤害，如果通过这种方式另一个人可以免受淫欲暴行的侵害。但是，正如我所说，这是另一个问题：目前让我们继续我们正在处理的问题：如果摆在我们面前的不可避免的条件是，要么说谎，要么遭受淫欲行为或某种可憎的玷污，即使说谎不伤害任何人，是否应该说谎。\par

26. 关于此事，若能先仔细讨论那些禁止说谎的神圣权威，则会有一些思考的空间：因为如果这些权威不留下余地，我们便徒然寻找漏洞；因为我们无论如何都必须遵守天主的命令，并在所有事上遵循天主的旨意，即使因遵守命令而受苦，我们也应当以平静的心接受；但若通过某种放宽，允许任何出口，那么在这种情况下，我们就不应拒绝说谎。圣经之所以不仅包含天主的命令，还包含义人的生平和品行，是为了：若偶然隐藏了我们应如何理解所命之事，可以通过义人的行为来理解。因此，除了那些可能具有寓意意义的行为——尽管无人怀疑它们确实发生过，正如旧约书中几乎所有的事件那样（因为谁敢断言那里有任何事不涉及预像的预言？），因为宗徒谈到亚巴郎的儿子们时（当然最容易说他们是按自然生育顺序出生并生活以繁衍百姓，因为没有生出怪物或异象来引导心灵走向某种预兆），然而他断言他们预像两个盟约；并且关于天主赐予祂的以色列子民以拯救他们脱离在埃及所受奴役的奇妙恩惠，以及在旅途中惩罚他们罪的惩罚，他说这些事发生在他们身上是作为预像——除了这些，在新约中由圣人们所行的事，那里有最明显的对品行的赞许以供我们效法，这些事可以作为理解圣经的例证，而这些事被编录在命令中。\par

27. 正如我们在福音中读到：\BibleQuote{你脸上挨了一记，把另一边也准备好。} 关于忍耐的榜样，没有比主自己的榜样更有力和卓越的了；但祂在被打脸时并没有说：看，这里是另一边，而是说：\BibleQuote{我若说得不对，请指证那不对之处；若说得对，你为什么打我？} 这里祂表明，准备另一边是在心中进行的。宗徒保禄也知道这一点，因为当他在公议会前被打脸时，他并没有说：打另一边；而是说：\BibleQuote{天主将要打你，你这粉白的墙！你坐在那里按法律审判我，却违反法律命令人打我？} 他以极深沉的洞察力看出犹太人的司祭职已变成这样：外表上洁净美丽，内里却充满污秽的贪欲；他说这些话时，在精神上看到那司祭职已准备通过主的惩罚而消逝；然而他早已准备好心，不仅接受脸上再挨打，而且为了真理忍受任何折磨，并且爱那些使他受苦的人。\par

28. 经上也记载：\BibleQuote{我却对你们说：总不可发誓。} 但宗徒自己在书信中却用了誓言。所以他表明如何理解所说的：\BibleQuote{我对你们说：总不可发誓：} 即，以免人因发誓而轻易发誓，从轻易变成习惯，再从习惯堕入背誓。因此我们发现他只在不说话的文字中发誓，因为文字中更有谨慎的预见，且没有不经意的舌头。而这的确出自邪恶，如经上说：\BibleQuote{此外任何事都是出于邪恶。} 但这并非出于他自己的邪恶，而是出于他们软弱中的邪恶，他即使用这种方式也在他们心中努力建立信心。因为他在说话时使用誓言，而不仅在写作时，我不知道圣经是否有关于他的记载。然而主说：\BibleQuote{总不可发誓：} 因为祂并没有给写作的人开此许可。然而，若宣称保禄违背了这条诫命，尤其是在为了万民的精神生活和救恩而书写并寄发的书信中，那便是亵渎；因此我们必须理解所写下的那个词\BibleQuote{总不可}，是为了说明：尽你所能，不要崇尚、不要喜爱，也不要仿佛为了某种好处而愉快地渴望发誓。\par

29. 正如：\BibleQuote{不要为明天忧虑，} 以及 \BibleQuote{所以不要忧虑吃什么、喝什么、穿什么。} 当我们看到主自己有一个钱袋，将人送给的东西放在里面，以便在必要时保留作为所需；并且宗徒们自己为弟兄们的贫困作了大量的准备，不仅为明天，甚至为即将来临的饥荒的更长时间，如我们在\BibleBook{宗徒大事录}{}中所读到的：这便足以清楚表明，这些诫命应这样理解：我们不应出于需要，因贪求暂时之物或害怕匮乏，而做任何属于我们工作的事。\par

30. 此外，宗徒们被告知，旅途中不可携带任何东西，而应依靠福音生活。主也在某处亲自指明祂说这话的原因，当祂加上说：\BibleQuote{工人配得他的工价：}祂充分表明这是许可，而非命令；免得有人若这样做，即在此宣讲圣言的工作中，从他所宣讲的对象那里收取任何生活所需，便以为自己做了什么不合法的事。然而，宗徒保禄充分证明，不做此事更为可赞：他虽曾说：\BibleQuote{凡在圣言上受教的，应与施教者分享一切财物，}并在许多地方表明，他所宣讲圣言的对象如此做是有益的，但他又说：\BibleQuote{然而，}他说，\BibleQuote{我没有使用过这权柄。}因此，主说这些话时，是赐予权柄，而非以命令束缚人。所以，一般而言，凡我们无法在言语中理解的，我们从圣人的行动中领悟当如何领会，否则易于偏向另一边，除非被榜样召回。\par

31. 因此，经上所写的：\BibleQuote{说谎的口杀害灵魂。}它所指的是哪一张口，是一个问题。因为当圣经一般提到口时，它指我们心中意念的所在，凡我们凭着真理说话时所发出的声音，都在那里被认可和决定。所以，凡认可谎言的，便是用心说谎。但是，一个人可能不会用心说谎，如果他口中说出了与心中所想不同的话，并且他知道这样做是为了避免更大的恶而承认一个恶，同时既不赞同这个恶也不赞同那个恶。持此说法的人认为，经上所写的\BibleQuote{心中说实话的人}也应如此理解。因为心中必须常说真理，但身体的口却不必总是如此，如果有任何为了避免更大的恶的理由，要求口中说出与心中所想不同的话。而确实存在心之口，这一点甚至可以从以下事实理解：哪里有言语，那里就有口，这并不荒谬。否则，如果说\Quote{谁在心里说话}就不恰当，除非认为心里也有一张口。虽然就在那处所写的\BibleQuote{说谎的口杀害灵魂}中，若考虑课文的上下文，或许可以理解为只指心之口。因为那里有一个隐晦的回应，向人隐藏，除非身体的口与其共鸣，否则心之口无法被人听闻。但圣经在该处说，那张口能到达充满全地的上主之神的听觉；同时在该处提到了嘴唇、声音和舌头；然而，感觉不允许将这些理解为关于心的，因为论到主说，所说的一切都瞒不过祂。而现在，以那种传到我们耳朵的声音所说的话，也同样瞒不过人。诚然，如此写着：\BibleQuote{智慧之神是慈爱的，必不豁免说邪恶之言的嘴唇：因为天主是祂肾脏的见证，是祂心中的真实探寻者，是祂舌头的听众。因为上主的神充满了全地，那包含万有的认识声音。因此，说非义之事的人不能隐藏；但审判惩罚时也必不放过他。因为在恶人的思想中要受审问，他的话声也要从上主那里达到，以惩罚他的不义。因为嫉妒的耳朵听闻一切，抱怨的喧嚣不能隐藏。所以，你们要戒绝无益的抱怨，制止舌头不毁谤人：因为隐晦的回应不会落空。但说谎的口杀害灵魂。}这样看来，它似乎是在警告那些认为他们所思所想是模糊和秘密的人。而这一点，它表明，在天主的耳中是如此清晰，以至于它甚至称之为\Quote{喧嚣}。\par

32. 同样地，我们在福音中明显找到心之口：以至于在某处，主被发现同时提及身体的口和心的口，祂说：\BibleQuote{你们还不明白吗？你们不晓得凡入口的，是进到肚腹，又落到茅厕里吗？惟独出口的，是从心里发出来的，这才污秽人。因为从心里发出来的，有恶念、凶杀、奸淫、邪淫、偷盗、妄证、亵渎：这些都是污秽人的。}此处，你若只理解为一张口，即身体的口，你将如何理解\InnerQuote{那些从口出来的，是从心里发出来的}？因为吐痰和呕吐也从口而出。除非一个人只在吃污秽之物时被玷污，而吐出时反而被玷污。但这若极其荒谬，那么我们就必须理解主是在解释心之口，当祂说\InnerQuote{那些从口出来的，是从心里发出来的}。因为偷窃可以而且经常在身体声音和口保持沉默的情况下进行；如果有人如此理解，即认为一个人只有在承认或揭露时才被偷窃之罪污染，而当他犯罪并保持沉默时却认为他未受玷污，那他一定是疯了。但是，实际上，如果我们把这句话指向心之口，那么任何罪都不能暗中犯下：因为罪除非从内在的那张口发出，否则不会犯。\par

33. 然而，就如同有人问圣经是在说什么口时说到：\Quote{说谎的口杀害灵魂。}同样也可以问，是对什么谎言的。因为它似乎特别指那由诋毁构成的谎言。圣经说：\Quote{要远离那无益的怨言，制止你们的舌头不事诋毁。}而这诋毁是出于恶意，当一个人不仅用口和身体的声音说出他所捏造的毁谤别人的话，甚至在不说话时也希望别人被看作如此；这实在是以心口诋毁；这事，圣经说，不能瞒过天主而不被知晓。\par

34. 因为在另一处写着的：\Quote{不要任意说谎。}他们说，这并不足以证明一个人不可使用任何谎言。因此，当一人说，按照这圣经的见证，我们必须如此憎恶各种各类的谎言，甚至人若想撒谎，虽然他没有撒谎，那愿望本身也应受谴责；且按此义解释，经文并非说：\Quote{不要使用任何谎言}，而是说：\Quote{不要任意使用谎言。}故人不可胆敢讲任何谎言，甚至不可愿讲任何谎言。另一人说：\Quote{不然，既然它说\InnerQuote{不要任意使用谎言}，是愿意我们从心口中根除并远离说谎：这样，虽然有些谎言我们须以肉口禁戒（如那些主要涉及宗教教义的），但有些谎言我们不必以肉口禁戒，若为了避免更大的恶而有理由；然而，我们须以心口完全戒除每一个谎言。}在此，必须理解\Quote{不要任意}的含义：即，这意志本身被视为心口，因此当为避免更大的恶而不得已说谎时，这不涉及心口。还有第三种理解，你可以如此理解\Quote{并非一切}这个词，即，除了某些谎言外，它允许你说谎。如同若说：\Quote{不要相信每一个人}，他并非主张谁也不应相信，而是说不要相信所有人，但有些人仍应相信。接着的话：\Quote{因为惯于说谎不会有利于善。}听起来似乎不是禁止说谎，而是禁止惯于说谎，即说谎的习性和喜爱。那必定会滑向这种习性的人，或是认为一切谎言都可随意使用（如此他连在虔诚和宗教教义中犯的谎言也不回避；在所有谎言中，甚至所有罪恶中，你能轻易找到比这更可恶的事吗？）或是将意志倾向某种谎言（无论多轻易、多无害），以至于他不是为了逃避更大的恶而不得已说谎，而是甘心情愿地喜欢说谎。所以，在这句话中可以理解三种意思：要么\Quote{不仅不要说任何谎言，甚至连愿说也不可}；要么\Quote{不要愿意，但当必须避免更坏的事时，不得已地说谎}；要么\Quote{并非一切谎言}，意即除了某些谎言外，其余允许。这三种中，一种支持那些主张永不可说谎的人，两种支持那些认为有时可以说谎的人。但接下来的\Quote{因为惯于说谎不会有利于善}，我不知道是否支持这三种中的第一种；除非如此：既然为成全者不仅不可说谎，甚至不可愿说谎，但连初学者也不可惯于说谎。也就是说，当设下规则无论如何不仅不可说谎，甚至连愿说谎也不行，而这一规则受到例子的反驳（因为有些谎言甚至得到大权威的认可），那么可以反驳说：那些确实是初学者的谎言，在此生中有某种怜悯的责任；但所有谎言都是恶的，应被成全和属灵的人完全避免，以至于连初学者也不被允许惯于说谎。因为我们已论及埃及的接生婆，她们在说谎时被赞许，是就着成长和进步到更好之事的应许而言；因为当一人出于怜悯而为了救任何人的性命（即使是必死的生命）而说谎时，这是迈向爱真实和永恒的灵魂得救的一步。\par

35. 此外，所写的：\Quote{凡说谎者你都要毁灭。}一人说这里没有例外，一切谎言都被谴责。另一人说：确实如此，但这是指从心里说谎的人，如我们前面所争论的；因为那在心里说真理的人，是厌恶说谎的必要性，将此视为道德生活的惩罚。另一人说：确实，天主将毁灭所有说谎的人，但不是所有谎言；因为有一种谎言是先知当时所暗示的，其中无人被宽恕；那就是，若一个人拒绝承认自己的罪，反倒辩护，且不愿做补赎，以致不满足于行不义，还希望被视为义人，不屈服于告解的良药。正如这些词本身的区别似乎暗示的：\Quote{你恼恨所有作恶的人}；但若他们悔改，在告解中说出真理，你就不毁灭他们，好使他们因行真理而来到光明中，正如在\BibleBook{若望}{福音}中所说的：\BibleQuote{但行真理的，却来就光明。}\par

36. 關於假見證——這記在法律十誡中——確實絕不可爭辯說，內心可以保持對真理的愛，同時卻向那見證所指向的人作出假見證。因為，當只對天主說時，真理只應在內心懷抱；但當對人說時，我們也必須用肉體的口說出真理，因為人不是內心的鑒察者。而關於見證本身，合理地問：人是誰的見證？因為我們說話的對象並非都是我們做見證的對象，只有那些藉著我們來得知或相信真理是適宜且應當的人，才是我們做見證的對象；例如法官，使他判決不犯錯；或接受宗教學說者，使他在信仰上不犯錯，或不受教師權威的影響而動搖。但當詢問你或想知道某事的人，尋求與他無關或不適宜他知道的事時，他要求的不是見證，而是告密者。因此，你若向他說謊，或許可以免於假見證，但絕不能免於謊言。所以，保留\BibleQuote{決不可作假見證}這一條件，問題在於：有時說謊是否合法？或者，如果任何說謊都是假見證，那麼就要看它是否可允許有補償，即是說，為了避免更大的罪而說謊；就如經上所寫的：\BibleQuote{應孝敬父親和母親，} 但當更優先的義務要求時，這條誡命可被忽略；因此，那被主親自召喚去宣講天主之國的人，被禁止為父親舉行最後的葬禮。\par

37. 同樣，關於所寫的：\BibleQuote{領受聖言的兒子必遠離毀滅；但他領受，是為自己領受，他口中沒有謊言。} 有人可能說，這裡所記的\BibleQuote{領受聖言的兒子}應理解為天主聖言，即真理。因此，領受真理的兒子必遠離毀滅指的是所寫的：\BibleQuote{祢必消滅一切說謊言的人。} 但當接著說\BibleQuote{他領受，是為自己領受}時，這暗示了什麼，不就是宗徒所說的：\BibleQuote{各人應考驗自己的工作，這樣，他只在自身內誇耀，而不在他人身上誇耀嗎？} 因為那領受聖言，即真理，不是為自己，而是為討人喜歡的人，當他看到別人能被謊言討好時，就不會保持真理。但誰是為自己領受的，他口中就沒有謊言：因為即使取悅人的方式是說謊，那領受真理不是為取悅人而是為取悅天主的人，就是為自己領受真理。因此，沒有理由說這裡\BibleQuote{祂必消滅一切說謊言的人}但並非所有謊言：因為所有謊言，普遍地，都被這句話砍斷：\BibleQuote{他口中沒有謊言。} 但另有人說，應該像宗徒保祿理解主的話那樣來理解：\BibleQuote{但我對你們說：總不可起誓。} 因為這裡也砍斷了一切起誓；但從內心之口而言，起誓絕不應出於意志的認可，而是出於他人軟弱的需要；也就是說，出於他人的\Quote{惡}，當顯示除非藉著誓言，否則無法使他人相信所說的話時；或出於我們自己的\Quote{惡}，即當我們仍被這必死生命的皮膚包裹著，無法顯露我們的心時；如果我們能夠做到，就無需起誓了。此外，在整段話中，如果\BibleQuote{領受聖言的兒子必遠離毀滅}這句話指的是那藉著祂萬物受造、永恆不變的真理本身；那麼，由於宗教學說致力於引導人默觀這真理，似乎可以認為\BibleQuote{他口中沒有謊言}這句話的意思是他不說任何與道理有關的謊言。這種謊言無論如何不可因補償而說，必須徹底、首先避免。或者，如果\BibleQuote{沒有謊言}這句話若不指向一切謊言就顯得不合理，那麼他口中就應像前面所論證的那樣，理解為內心之口，這是持有時可以說謊觀點的人的意見。\par

38. 确然无疑，尽管争论双方各执一词，有人坚称绝不可说谎，并为此引述神圣的见证；另有人反驳，甚至在神圣见证的字里行间为谎言寻找立足之地；然而无人能说，他从圣经的榜样或话语中发现，任何谎言应被视为可爱之物，或不应被憎恶；不过有时藉着说谎，你必须做你所憎恶的事，以免那更应被憎恶的事发生。但正是在这里人们犯错了：他们将珍贵置于卑下之下。因为当你承认某种邪恶可以接受，以避免另一个更严重的邪恶时，每个人不是依据真理的准则，而是依据自己的贪欲和习惯来衡量邪恶，认为他自己更惧怕的那个更严重，而非实际上更应逃避的那个。这一切过犯皆由爱德之颠倒而生。因为我们有两种生命：一种是永恒的，是天主所应许的；另一种是暂时的，我们现在就在其中。当一个人开始爱这暂时的生命胜过那永恒的生命时，为了他所爱的，他认为万事皆可做；在他眼中，没有比伤害今世生命、或以不义不法的方式剥夺其任何利益、或施加死亡将其完全夺去更严重的罪了。因此，盗贼、土匪、暴徒、酷刑者、杀人者，在他们眼中，比好色、酗酒、奢侈之人更可恨，倘若这些人不打扰他人。因为他们不明白，也毫不关心，这些人得罪了天主；并非对天主有任何不便，而是自招有害的祸患；因为他们败坏祂赐予他们的恩赐，即暂时的恩赐，并因自己的败坏而远离永恒的恩赐；尤其是，如果他们已经开始成为天主的圣殿；关于这一点，宗徒对所有基督徒这样说：\BibleQuote{你们不知道你们是天主的圣殿，天主的圣神住在你们内吗？谁若毁坏天主的圣殿，天主必要毁坏他，因为天主的圣殿是圣的，这圣殿就是你们。}\BibleRef{格前 3:16-17}\par

39. 确实，所有这些罪，无论是那些损害他人今世安适的，还是那些自甘堕落而不违背他人意愿的：所有这些罪，即使它们似乎对今世生命有益，以获取任何愉悦或利益（因为无人以其他目的和结局犯这些罪），但就那永永远远的生命而言，它们确实缠累并完全阻碍人。但其中有些只阻碍行恶者本身，另有些则同样阻碍那些承受恶行的人。因为就人们为了今世利益而保存的事物而言，当这些事物被行恶者夺走时，只有如此行的人犯罪并被阻碍进入永生，而非受其害的人。因此，即使一个人同意他人从他那里取走这些东西，或为了避免行恶，或为了避免在这些事情上遭受更大的不便，他不仅没有犯罪，而且在一种情况下他行动英勇、值得称赞，在另一种情况下他有益且无可指责。但就那些为了圣德和宗教而保存的事物而言，当行恶者想要侵犯它们时，如果有条件提出且有机会，用较小的罪来救赎它们是合理的，但不能通过损害他人来做到。这样一来，那些为了避免更大的罪而采取的事，从此就不再是罪。因为就如对有用的事物，例如金钱或其他任何身体之利，为了更大的利益而舍弃的不称为损失；同样，对神圣的事物，为了避免更坏的事而接纳的，不称为罪。或者，若某人为了避免失去更多而放弃的被称为损失，那么这也可被称为罪；但为了避免更大的罪而采取此罪的必要性，无疑就如为了避免更大的损失而忍受较小的损失是正确的一样。\par

40. 为圣德而应当保持的事物如下：身体的贞洁、灵魂的贞洁和教义的真确。身体的贞洁，未经灵魂的同意与许可，无人能侵犯。因为凡违反我们意愿，未得我们授权，而由更大力量加于我们身体的事，并非淫乱。然而，许可可能有某种理由，但同意则毫无理由。因为我们同意时，便赞许并期望；但我们许可时，即使不情愿，也是由于要避免某种更大的卑污。真正对肉体淫乱的同意，也侵犯了心灵的贞洁。因为心灵的贞洁在于善的意愿和真诚的爱，除非我们爱慕并渴求那真理教导我们不应爱慕和渴求的事物，否则它不会败坏。因此，我们必须守护对天主和近人爱的真诚；因为在此，心灵的贞洁得以圣化。我们必须竭尽所能，并以虔诚的祈祷，尽力做到：当我们的身体贞洁遭受侵犯时，连那与肉体相缠绕的灵魂最外在的感官也不被任何愉悦所触动；但若不能如此，至少那不同意的心灵与思想，能保全其贞洁。至于我们在心灵的贞洁中应当守护的，就爱近人而言，是无害与仁爱；就爱天主而言，是虔敬。无害即不伤害任何人；仁爱即向所能及之人行善；虔敬即敬拜天主。至于宗教与虔敬的教义真确，则除非藉谎言，否则不会被侵犯；然而那至高的、内在的真理本身，即该教义所属者，绝不可能被侵犯。达到此真理，并完全存留于其中，彻底与它结合，只有在这\BibleQuote{必朽坏的必穿上不朽坏，这必死的必穿上不死的}时，才得允许。但是，因为今生的一切虔敬都是实践，我们借此趋向那生命；这实践从教义中获得引导，而教义藉人类言语和身体圣事的标记，暗示并传达真理本身。为此，那可能因说谎而受败坏的事物，最应保持不坏；如此，若心灵贞洁中有所受损，它仍可从中获得修复。因为教义的权威一旦败坏，便再无任何途径或方法可回归心灵的贞洁。\par

41. 由此得出以下结论：凡不损害虔敬教义、虔敬本身、无害与仁爱的谎言，可因身体贞洁的理由而被允许。然而，若有人立定心志如此爱慕真理——不仅包括沉思中的真理，也包括说出真实之事（即每一类事物中为真的那件事），并且除非心中所构想和洞见的，否则绝不从身体之口发出；以致他珍视诚实言说的美，不仅胜过金银、宝石和美好的田地，甚至胜过今生本身和身体的一切美善——我不知道有谁能明智地说此人错了。而且，若他以此事为优先，并将它珍视超过自己所拥有的一切，那么他也应正当地将它优先于他人的今世事物，因为他藉无害与仁爱有责任保护并帮助他人。因为他会爱完美的信德，不仅在于正当地相信那些由卓越而可信的权威向他所说的事，也在于忠实地说出他自己判断当说且将要说的话。因为\OriginalTerm{信}{fides}在拉丁语中得名于所说之事得以实现；由此可见，说谎时便没有展现信德。即使当一个人说谎时，其谎言不带来任何不便或有害的伤害，且其意图是保护自己的生命或身体的贞洁，此时信德所受破坏较小，但信德仍被破坏了，而信德是应当在心灵的贞洁与圣德中保持安全的事物。因此，我们被真理本身——那超越万有、独一不可战胜的真理——所迫，而非被多数陷于错误的人的意见所迫，甚至将完美的信德置于身体的纯洁之上。因为心灵的贞洁是秩序井然的爱，它不会将较大的置于较小的之下。凡身体中可被侵犯的，都小于心灵中可被侵犯的。的确，当一个人为身体的贞洁而说谎时，他确实看到自己的身体面临被败坏的危险——不是来自他自己的情欲，而是来自他人的情欲——但他谨慎，至少避免因许可而成为同谋。然而，这许可，除了在心灵中，又能在哪里呢？因此，甚至身体的贞洁也只能在心灵中被败坏；若心灵不同意也不许可，那么无论他人情欲对身体做了什么，都不能正当地说身体的贞洁被侵犯了。由此可见，心灵的贞洁更必须在心灵中得以保存，而身体贞洁的守护亦在其中。因此，就我们所能而言，两者都当以圣洁的品行和言谈围护起来，以免从别处被侵犯。但当两者不可兼得时，哪一个相比之下应被轻视？有谁看不出来呢？当他看到哪一个应被优先考虑——心灵相对于身体，或身体相对于心灵；以及哪一罪更应避免——是许可他人的行为，还是亲自犯罪。\par

42. 凡所讨论的，如今已清楚显明：圣经的那些见证并无别的含义，只在于我们绝不可在任何时候说谎；因为在圣人的品行和行动中，就那些不涉及比喻性含义的经书（诸如宗徒大事录中的历史）而言，找不到任何值得效法的说谎例子。至于我们主在福音中所说的一切话，在那些较为无知的人看来像是谎言，其实都是比喻性的含义。至于宗徒所说的：\Quote{对一切人，我就成为一切，为的是要赢得一切人；}正确的理解是：他这样做并非靠说谎，而是出于同情；他以极大的爱德对待他们，解救他们，仿佛他自己处于那恶中，他愿使人们从那里痊愈。因此，在虔敬的教义中绝不能说谎：这是一种极恶的恶行，是可憎之谎的第一类。绝不能说第二类谎，因为不可错待人。绝不能说第三类谎，因为我们不可为了谋取某人之利而损害他人。绝不能说第四类谎，即出于说谎之欲的谎，这本身就是邪恶的。绝不能说第五类谎，因为连真理本身也不可出于取悦人的目的而说出，何况谎言本身作为一种谎，是污秽的？绝不能说第六类谎，因为不可为了任何人的现世便利和安全而败坏见证之真理。但不可借助谎言引导任何人获得永恒的救恩。因为人不是因着那转变他之人的坏品行而被转变为好品行的：如果对他这样做是适宜的，那么他在转变后也应当对他人这样做；这样他转变的就不是好品行，而是坏品行，因为那转变他时所行的，被当作他转变后要效法的榜样。第七类谎也绝不能说，因为不可将任何人的利益或现世福祉置于信德的成全之上。即使有人因我们的正义行为而被如此恶意地触动，以致心智变得更糟，离虔敬更远，也不应因此放弃正义行为：因为我们必须首要持守的是，我们应当召唤并邀请那些我们如同爱自己一样去爱的人；且要以最勇敢的心志领受宗徒的那句名言：\Quote{对于一些人，我们是致死的死臭；对于另一些人，我们是生活的生香；谁又能胜任这些事呢？}第八类谎也绝不能说，因为在善事中，心灵的贞洁大于身体的贞洁；在恶事中，我们自己所做的恶大于我们所容忍的恶。然而，在这八类谎中，一个人说谎时，他越接近第八类，罪就越轻；越偏离向第一类，罪就越重。但若有人以为某种谎不是罪，他将严重自欺，因为他自认诚实，实则欺骗他人。\par

43. 此外，如此巨大的盲目占据了人的心思，以至于对他们而言，若我们宣称某些谎不是罪，这还不够；他们还必须宣称，若我们拒绝说谎，在某些事上这就是罪。他们因辩护说谎而到了如此地步，甚至宣称连所有谎中最可憎邪恶的第一类，宗徒保禄也曾使用过。因为在写给迦拉达人的书信中——这书信如同其他书信一样，是为宗教和虔敬的教义而写的——他们说保禄在谈到伯多禄和巴尔纳伯时说了谎：\Quote{我一看见他们不按福音的真理正直而行。}因为他们企图为伯多禄从错误和所陷入的邪路上辩护，却通过割裂和曲解圣经的权威，试图推翻那为人人带来救恩的宗教之路本身。他们并未看到，他们不仅指控宗徒在虔敬的教义中说了谎，还指控他发了虚誓——因为保禄在叙述那事之前说：\Quote{我所写给你们的，看，在天主面前，我并不说谎。}但现在是时候为这场辩论设定界限了；在考虑和处理此事时，最适宜首先牢记并祈祷的，正是同一宗徒所说的话：\Quote{天主是忠信的，他决不允许你们受试探超过你们所能承受的，而且随着试探，他也要开一条出路，叫你们能承受得住。}\par

\SourceRecord{Translated by H. Browne.；Nicene and Post-Nicene Fathers, First Series；Vol. 3.；Edited by Philip Schaff.；Buffalo, NY: Christian Literature Publishing Co.,；1887.；<http://www.newadvent.org/fathers/1312.htm>.}
